\chapter{Translation}

\section{Agreement at every lag}

For two views $a$ and $b$ of the same extents, the agreement at a lag $\ell$ is
\[
C(\ell) \;=\; \sum_{v} a(v)\, b(v + \ell),
\]
and the lag with the largest agreement is the translation between them. Computed lag by lag this
costs the volume times the number of lags. Computed through the convolution theorem it costs one
transform of each view, one pointwise product, and one inverse transform, for every lag at once.

\section{Phase correlation}

Phase correlation rests on the Fourier shift theorem. If two images differ by a displacement
$(\Delta x, \Delta y)$, their Fourier transforms are related by
\[
I_2(\omega_x, \omega_y) \;=\; I_1(\omega_x, \omega_y)\, e^{-j(\omega_x \Delta x + \omega_y \Delta y)},
\]
so the normalised cross power spectrum is a pure phase whose inverse transform is a delta at the
displacement. The method extends to $N$ dimensions directly. Under translation the motion model
for $N$ dimensional data is a rank one tensor, and a higher order singular value decomposition of
the phase correlation separates the displacement along each dimension independently.

Sub-pixel precision is recovered from the shape of the peak. The peak of phase correlation is a
Dirichlet kernel, closely approximated by a sinc, and fitting a sinc, a Gaussian or a triangle to it
recovers the fraction of a pixel. Foroosh, Zerubia and Berthod derive the sub-pixel extension in
closed form. An upsampled matrix multiplication of the discrete Fourier transform gives arbitrary
sub-pixel precision without transforming a larger grid, and is how the widely used
\texttt{register\_translation} of scikit-image works.

\section{Drift correction}

Time-lapse imaging uses exactly this to remove the drift of a subject in its view. In published
microscopy protocols phase correlation registers each time point of a four dimensional recording
while keeping the movements of the objects inside it measurable, and in long recordings the global
movement of the subject is estimated during acquisition and fed back to the stage so the subject
stays in the field of view. Video stabilisation does the same for a hand held camera. In every case
the translation of the whole view is removed before any object is tracked, which is the order the
first chapter states.

\section{Where translation stops}

A translation has as many parameters as the view has axes and explains only the motion every point
shares equally. A view that turns, zooms or changes perspective moves points by different amounts
in different places, and its residual after the best translation is a structured field, not
scatter. The next chapter takes those transforms in turn.
